\paragraph{Autonomous native repair from an unmodified strong starting route.}
The unified controller selects connections from each accepted state's native STA, queries capacity-feasible extensions, constructs a joint conflict scope when necessary, and lets the native serial router choose every replacement path. Occupancy is updated after each attachment. Complete candidates receive independent physical checks and native FULL/STA; rejected states leave the immutable accepted route unchanged. An exact state/request cache shares repeated attempts across the two policy branches. The initial R1 route is the sealed, unmodified strong route, not an injected intermediate repair.

Single attempts on 6587/pin~4 and 6590/pin~4 preserve the complete timing plateau and are rejected. Joint repair of these targets, with four automatically identified blocking branches, reaches 5.854659868020917 ns from 6.018101128546505 ns, improving 246 setup components without regressions. Both policies accept. Fresh STA selects 76/pin~1 and 48/pin~1; their joint successor reaches 5.832274663220005 ns. Strict CPD accepts this second step, whereas endpoint-Pareto rejects its 0.371651154 ns loss at \texttt{Mp[12].D[0]} and retains its own first-step state. The strict branch's next six-connection frontier produces a legal 6.285112431 ns candidate, which is rejected. Thus this fully automatic native trajectory obtains the earlier best local CPD, 3.087792\% below the unmodified strong R1, without an injected first repair. It uses 40 SB pairs, 970 CB sources and planar wirelength 218245, compared with the original 39, 969 and 218164. This is one-circuit development evidence, not general promotion.

Relative to the original R1 endpoint vector, strict CPD improves 254 setup components and regresses one; endpoint-Pareto improves 246 and regresses none. The fixed initial 0.2 ns endpoint window shrinks from three components to one and two respectively. Windows recentered on each final state's own worst slack contain eighteen and fourteen components; these moving counts must not be read as the same fixed cohort. The unchanged native default constraints contain \texttt{virtual\_io\_clock}, \texttt{clk} and \texttt{clk\_mem}; no new SDC is supplied. The fixed vector contains 14032 setup/domain components over 13744 endpoints. CPD and vector decisions use the native full-precision numbers with no acceptance epsilon, while increased negative hold violations remain forbidden.

\paragraph{Released scope contributes beyond selection and full-device search.}
A matched native control starts from the same unmodified R1 and receives the same current-STA selector, full-device search, native cost objective and single/joint acceptance opportunities, but releases only the target branches. Both policy branches remain at 6.018101128546505 ns. The successful first conflict-scope repair additionally releases 1429/pin~1, 1967/pin~1, 6563/pin~1 and 7963/pin~2. Therefore selection and extra full-device target repair alone do not explain this first gain. Scope differs intentionally in this ablation; the earlier native/external engine comparisons retain their matched scopes. Exact previously completed single attempts and joint query results are reused after source state, request, normalized native command, objective and input checks. A preliminary reuse-controller assertion incorrectly included batch registration provenance in input identity; it stopped before any native call. Its failure is retained, and an independent corrected controller excludes only registration metadata not consumed by the VPR command.

\paragraph{Complete prototype cost is dominated by state preparation and checking.}
The two-policy conflict-scope R1 controller takes 631.358 s, including rejected work and cache handling. Nonoverlapping child times are 249.654 s for query/candidate audits, 196.460 s for candidate FULL/STA wrappers, 73.792 s for capacity/relaxed-query preparation and joint search, 72.494 s for native repair wrappers, and 23.328 s for scope preparation; 15.631 s occurs outside children. The native transactions themselves total only 1.979 s and are nested within the native wrappers. Five distinct native transactions make 66 connection-search attempts; the endpoint branch reuses common-state work. The target-only control costs 102.544 s incrementally, with one new native transaction and five searches; its reused work prevents interpreting the ratio as independent policy speedup. Continuing the earlier accepted native state costs 177.302 s, including a rejected candidate, versus 1.406 s inside its transaction. These measurements begin from existing legal routes and exclude generating those original routes. They are additional prototype-flow costs, not formal cold-start complete-routing benchmarks, and use no historical U timing.

\paragraph{Protecting a previous victim moves the bottleneck; joint capacity constrains the remaining potential.}
At 5.832274663 ns, only 5490/pin~5 and 6591/pin~4 among the six worst connections have positive relaxed fixed-tree delay potential, approximately 7.639 ps each. The first joint scope releases 21 connections and worsens 1388 setup components. A current-STA rule identifies the released 6589/pin~4 prefix as having only about 7.639 ps headroom. Reserving this prefix still admits compatible improving target paths, but their new conflict scope affects 25 connections. Native recovery preserves the protected prefix and lowers both target delays, yet the critical loss migrates to 6592/pin~4. The candidate remains legal at 6.285112430504114 ns, worsens 1387 setup components and is rejected by both policies. This control takes 87.245 s, including 1.757 s inside native recovery and 47 connection-search attempts. Target potential and a legal transaction do not imply CPD improvement.

The next diagnostic protects all eight current non-target prefixes whose slack headroom is within the maximum target gain, including newly encountered competitors. Each target separately retains its 7.639 ps potential, but joint conflict splitting identifies capacity-one IPIN 82116 as necessary for either strict local improvement. Forbidding this IPIN on either target produces an exact distance of 0.856978293883870 ns, its original delay, so the two gains cannot coexist with these eight prefixes fixed. This is a conditional residual-graph result, not a global routing or endpoint optimum. Protection is stronger than a timing budget: equal-delay victim relocation, different retained-tree scope and upstream changes remain untested possibilities. The next registered branch permits current-STA-budgeted victim alternatives and compares a slack-offset objective with raw connection-delay selection, retaining whole-transaction FULL/STA acceptance.
